
This work tested whether fixed statistical thresholds derived from ImageNet literature could enable automated semantic versioning (MAJOR/MINOR/PATCH classification) for machine learning datasets. Experiments on 14 dataset pairs falsified this hypothesis: the classifier achieved 28.57\% accuracy (vs. 85\% target), with 25\% precision and 25\% recall for MAJOR change detection (vs. 70\% and 85\% targets).

The primary failure mode was dual breakdown: the classifier missed 3 out of 4 true MAJOR changes (75\% false negative rate) while incorrectly flagging 6 non-MAJOR changes as MAJOR (86\% false positive rate among MAJOR predictions). Root cause analysis identified 18.8× variance in drift scores across datasets with similar ground truth labels, demonstrating that drift magnitude is dataset-relative rather than absolute.

These results provide quantitative evidence that ImageNet-derived thresholds do not generalize to the tested dataset collection without substantial recalibration. The below-random accuracy (28.57\% vs. 33\% baseline) indicates that fixed thresholds provide worse than uninformative signal for semantic version classification in this evaluation.

\textbf{Limitations:} Ground truth labels were literature-derived rather than performance-validated. Dataset coverage included primarily NLP benchmarks (13 of 14 valid examples), with limited vision evaluation. The MNIST result represents invalid cross-dataset comparison. These limitations bound the scope of claims but do not change the negative result given the substantial performance gaps observed.

\textbf{Scientific contribution:} This work provides empirical falsification of fixed-threshold semantic versioning across multiple dataset types, with quantitative measurement of failure magnitude (-56.43pp accuracy gap, -45pp precision gap, -60pp recall gap). While negative results are less recognized than positive findings, rigorous falsification eliminates unproductive research directions. The documented 18.8× drift variance and systematic failure across all severity levels provide evidence for pursuing adaptive, supervised, or performance-based approaches rather than universal fixed thresholds.

Future work should validate performance-based ground truth, expand vision dataset coverage through manual downloads, and evaluate adaptive calibration methods that establish dataset-specific thresholds from observed version transitions.
